\chapter{Introduction}
\label{ch:introduction}

Artificial intelligence has grown explosively in recent years, not only as
a field of scientific research but as a technology that is rapidly entering
everyday life. The trend is unmistakable, and we may well be living through
one of the most significant shifts in the history of technology.

The story of machine learning reaches back to the 1940s and 1950s, with the
first attempts to build computational models inspired by the brain. The
perceptron, proposed by Frank Rosenblatt in 1958~\citep{rosenblatt1958perceptron},
was one of the earliest learning algorithms able to classify linearly
separable patterns. Its theoretical limitations, made precise by Minsky and
Papert~\citep{minsky1969perceptrons}, contributed to the first ``AI winter'':
a prolonged period of skepticism and reduced funding. The revival came in the
1980s with the backpropagation algorithm~\citep{rumelhart1986backprop}, which
made it possible to train multilayer networks efficiently. Even so, vanishing
gradients and the limited computing power of the time kept these techniques in
the background for years.

The transformation we now call \emph{deep learning} arrived in 2012, when
AlexNet~\citep{krizhevsky2012alexnet} won the ImageNet competition and
slashed the classification error of traditional methods. This success rested
on the convergence of three factors: large labelled datasets, accessible
specialized hardware (GPUs) capable of training models with millions of
parameters, and training techniques that overcame the historical obstacles of
deep networks. Since then, deep learning has proven remarkably effective at
image recognition, natural language processing, content generation, and
protein structure prediction, among many other tasks.

The architectures behind these achievements --- convolutional, recurrent, and
attention-based networks --- were conceived to operate on data laid out in
regular Euclidean domains: grids of pixels, temporal sequences, tables of
features. As the community began applying them to data with intrinsically
richer structure (molecules, social networks, three-dimensional surfaces), it
became clear that the assumption of Euclidean regularity was a fundamental
limitation --- conceptual, not merely technical. This realization gave rise to
a new paradigm, \emph{geometric deep learning} (GDL), which inverts the
traditional relationship between data and architecture: instead of forcing the
data to fit pre-existing models, the geometric properties of the data dictate
the structure of the model.

\section{What is geometric deep learning?}
\label{sec:what-is-gdl}

Traditional architectures are designed for data with regular Euclidean
structure. Convolutional networks exploit the regular grid of an image, where
every pixel has a fixed number of neighbours in predetermined positions;
recurrent and attention models process sequences endowed with a natural
temporal order. This dependence on regular domains breaks down when the data
does not fit such templates.

Consider analysing a social network: users and their connections form an
irregular graph in which each node may have an arbitrary number of neighbours,
with no natural ordering or grid. Applying a convolutional network directly
would require flattening the graph into a Euclidean representation (such as an
adjacency matrix), discarding crucial structural information and introducing a
flood of zeros that make the computation wasteful. The analysis of molecules
in drug discovery is equally telling: molecules are naturally graphs whose
nodes are atoms and whose edges are chemical bonds, and their pharmacological
behaviour is governed by three-dimensional geometry and bond topology.
Classical methods relied on hand-crafted descriptors or string encodings (such
as SMILES), losing geometric information essential to biological activity. The
same ``curse of non-Euclideanity'' appears in the analysis of 3D shapes, where
surfaces admit no natural parametrization on a regular grid and techniques such
as voxelization or multi-view projection introduce artefacts and erase fine
geometric detail.

\Cref{fig:data-types} illustrates the three fundamental categories of data
according to their geometric structure: Euclidean data with a regular grid,
graphs with irregular topology, and manifolds with continuous geometry.

\begin{figure}[h]
\centering
\begin{tikzpicture}[scale=0.9]
% --- Left panel: Euclidean (regular grid) ---
\begin{scope}[shift={(-5,0)}]
    \begin{scope}[on background layer]
        \fill[gdlblue!15, rounded corners=2pt] (-0.3,-0.3) rectangle (3.3,3.3);
    \end{scope}
    \foreach \x in {0,...,3} \foreach \y in {0,...,4}
        \draw[gdlgray!50] (\x*0.75, \y*0.75) -- (\x*0.75+0.75, \y*0.75);
    \foreach \x in {0,...,4} \foreach \y in {0,...,3}
        \draw[gdlgray!50] (\x*0.75, \y*0.75) -- (\x*0.75, \y*0.75+0.75);
    \foreach \x in {0,...,4} \foreach \y in {0,...,4}
        \fill[gdlblue!60] (\x*0.75, \y*0.75) circle (2.5pt);
    \node[below, font=\bfseries, text=gdlblue!70!black] at (1.5, -0.7) {Euclidean};
\end{scope}
% --- Centre panel: Graph (irregular topology) ---
\begin{scope}[shift={(0.5,0)}]
    \coordinate (n1) at (0.3, 2.8); \coordinate (n2) at (1.8, 3.0);
    \coordinate (n3) at (0.0, 1.5); \coordinate (n4) at (1.5, 1.6);
    \coordinate (n5) at (2.8, 2.0); \coordinate (n6) at (0.5, 0.2);
    \coordinate (n7) at (2.2, 0.3);
    \foreach \a/\b in {n1/n2,n1/n3,n1/n4,n2/n4,n2/n5,n3/n4,n3/n6,n4/n5,n4/n6,n4/n7,n5/n7,n6/n7}
        \draw[gdlgreen!50, thick] (\a) -- (\b);
    \foreach \n in {n1,n2,n3,n4,n5,n6,n7} \fill[gdlgreen!70] (\n) circle (4pt);
    \node[below, font=\bfseries, text=gdlgreen!70!black] at (1.5, -0.7) {Graph};
\end{scope}
% --- Right panel: Manifold (continuous geometry) ---
\begin{scope}[shift={(5.5,0)}]
    \shade[ball color=gdlorange!25, opacity=0.9]
        (0.0, 1.5) .. controls (-0.5, 2.8) and (1.0, 3.5) .. (1.8, 3.0)
        .. controls (2.6, 2.5) and (3.3, 2.8) .. (3.2, 1.8)
        .. controls (3.1, 0.8) and (2.5, -0.2) .. (1.5, 0.0)
        .. controls (0.5, 0.2) and (0.5, 0.3) .. (0.0, 1.5);
    \draw[thick, gdlorange!60]
        (0.0, 1.5) .. controls (-0.5, 2.8) and (1.0, 3.5) .. (1.8, 3.0)
        .. controls (2.6, 2.5) and (3.3, 2.8) .. (3.2, 1.8)
        .. controls (3.1, 0.8) and (2.5, -0.2) .. (1.5, 0.0)
        .. controls (0.5, 0.2) and (0.5, 0.3) .. (0.0, 1.5);
    \draw[gdlorange!30, thin] (0.5, 1.2) .. controls (1.2, 1.8) .. (2.5, 1.5);
    \draw[gdlorange!30, thin] (0.8, 2.2) .. controls (1.5, 2.6) .. (2.6, 2.2);
    \node[below, font=\bfseries, text=gdlorange!70!black] at (1.5, -0.7) {Manifold};
\end{scope}
\end{tikzpicture}
\caption[Types of geometric data]{The three fundamental kinds of data by
geometric structure: Euclidean data on a regular grid (left), graphs with
irregular topology where each node may have a different number of neighbours
(centre), and manifolds with continuous geometry requiring local
parametrizations (right).}
\label{fig:data-types}
\end{figure}

Geometric deep learning is a mathematically principled response to these
limitations. Rather than adapting the data to existing architectures --- a
process that invariably loses information --- it designs architectures that
respect and exploit the intrinsic geometric structure of the data. The central
observation is that many real-world domains carry natural symmetries and
geometric structure that can be exploited to improve both computational
efficiency and predictive performance. By building these symmetries directly
into the network, GDL not only resolves the representational problems above but
also guarantees desirable properties such as equivariance, robust
generalization, and parameter efficiency.

The full mathematical foundations --- group theory, the notions of invariance
and equivariance, and the unifying ``5 Gs'' framework --- are developed in
\Cref{ch:foundations,ch:priors}. For now it is enough to understand that GDL
provides a theoretical framework unifying seemingly disparate architectures
--- convolutional networks, graph networks, and attention mechanisms --- under
a common geometric paradigm.

\section{Symmetry as a design principle}
\label{sec:symmetry-principle}

The organizing idea of geometric deep learning is \emph{symmetry}. A symmetry
of a problem is a transformation of the data that leaves the quantity of
interest unchanged. Rotating an image of a cat still shows a cat; relabelling
the nodes of a molecular graph does not change the molecule. Such
transformations form a \emph{group}, and demanding that a model behave
consistently under that group is a powerful constraint on its architecture.

Two notions make this precise. A function is \emph{invariant} when its output
does not change as the input is transformed by the group --- appropriate, for
example, for whole-image classification. It is \emph{equivariant} when
transforming the input produces a correspondingly transformed output: if
$f$ is equivariant to a group $\group$ acting on its input and output spaces,
then
\[
  f(g \cdot x) = g \cdot f(x) \qquad \text{for all } g \in \group .
\]
Equivariance is the property we want from the intermediate layers of a network:
each layer should transform predictably under the symmetry so that structure is
preserved as information flows through the model, with an invariant readout only
at the end. Because equivariant maps compose into equivariant maps, this
principle yields a constructive recipe for entire architectures rather than a
property to be checked after the fact. These ideas are formalized in
\Cref{ch:priors}.

\section{The 5 Gs: a roadmap}
\label{sec:five-gs}

Geometric deep learning organizes its domains and symmetries around five
recurring themes --- the ``5 Gs'' that give this book its subtitle: grids,
groups, graphs, geodesics, and gauges~\citep{bronstein2021geometric}.
\Cref{tab:five-gs} summarizes them and points to the chapter where each is
developed.

\begin{table}[h]
\centering
\caption{The five Gs of geometric deep learning.}
\label{tab:five-gs}
\begin{tabular}{@{}lp{6.2cm}p{3.6cm}@{}}
\toprule
\textbf{Theme} & \textbf{Idea} & \textbf{Chapter} \\
\midrule
Grids & Regular Euclidean domains with translation symmetry & Convolutional networks (\Cref{ch:grids}) \\
\addlinespace
Groups & Global symmetries beyond translation, such as rotations and reflections & Group-equivariant networks (\Cref{ch:groups}) \\
\addlinespace
Graphs & Irregular relational domains with permutation symmetry & Graph neural networks (\Cref{ch:graphs}) \\
\addlinespace
Geodesics & Intrinsic geometry of manifolds: distances, curvature, and information flow & Learning on manifolds (\Cref{ch:geodesics}) \\
\addlinespace
Gauges & Local reference frames and the freedom in choosing coordinates & Gauge-equivariant networks (\Cref{ch:gauges}) \\
\bottomrule
\end{tabular}
\end{table}

These themes are not isolated. Grids are the most rigid domain and the
historical starting point; relaxing the requirement of a global grid while
keeping global symmetries leads to group-equivariant networks; discarding the
ambient grid altogether and keeping only relational structure leads to graphs;
endowing a domain with intrinsic distances and curvature leads to manifolds and
geodesics; and acknowledging that on a curved domain there is no canonical
choice of local frame leads to gauges. A single thread --- message passing ---
runs through all of them and exposes convolutional networks, graph networks,
and transformers as instances of one geometric construction.

\section{How to read this book}
\label{sec:how-to-read}

The book follows a logical progression from foundations to advanced models.
\Cref{fig:book-structure} shows the dependencies between chapters.

\begin{figure}[h]
\centering
\begin{tikzpicture}[
    node distance=0.9cm and 1.4cm,
    box/.style={rectangle, draw, rounded corners, minimum width=2.7cm,
        minimum height=0.8cm, align=center, font=\small},
    found/.style={box, fill=gdlblue!15},
    five/.style={box, fill=gdlgreen!18},
    prac/.style={box, fill=gdlpurple!16},
    arrow/.style={->, thick, >=stealth}
]
\node[found] (intro) {Ch.\ 1\\Introduction};
\node[found, below=of intro] (found) {Ch.\ 2\\Foundations};
\node[found, below=of found] (priors) {Ch.\ 3\\Geometric priors};
\node[five, below left=1.1cm and 0.2cm of priors] (grids) {Ch.\ 4--5\\Grids, Groups};
\node[five, below right=1.1cm and 0.2cm of priors] (graphs) {Ch.\ 6--8\\Graphs, Geodesics, Gauges};
\node[prac, below=2.6cm of priors] (apps) {Ch.\ 9--10\\Applications, Conclusion};
\draw[arrow] (intro) -- (found);
\draw[arrow] (found) -- (priors);
\draw[arrow] (priors) -- (grids);
\draw[arrow] (priors) -- (graphs);
\draw[arrow] (grids) -- (apps);
\draw[arrow] (graphs) -- (apps);
\node[right=1.6cm of priors, font=\footnotesize, align=left] {
    \tikz\fill[gdlblue!15] (0,0) rectangle (0.3,0.3); ~Foundations\\[2pt]
    \tikz\fill[gdlgreen!18] (0,0) rectangle (0.3,0.3); ~The Five Gs\\[2pt]
    \tikz\fill[gdlpurple!16] (0,0) rectangle (0.3,0.3); ~Practice};
\end{tikzpicture}
\caption{Structure of the book and dependencies between chapters.}
\label{fig:book-structure}
\end{figure}

\textbf{Part I (Foundations).} \Cref{ch:foundations} establishes the
mathematical language --- groups, representations, homogeneous spaces, and
manifolds. \Cref{ch:priors} introduces invariance and equivariance and the
blueprint that turns a symmetry into an architecture.

\textbf{Part II (The Five Gs).} \Cref{ch:grids,ch:groups,ch:graphs,ch:geodesics,ch:gauges}
develop the five themes in turn, from convolutional networks on grids to
gauge-equivariant networks on manifolds, showing how each extends the previous
one and how message passing unifies them.

\textbf{Part III (Practice).} \Cref{ch:applications} surveys applications
across computer vision, the physical sciences, and relational data, and
\Cref{ch:conclusion} synthesizes the unified view and outlines open problems.

This is a theoretical book, with emphasis on the mathematical rigour of its
definitions, results, and the formal connections between the frameworks it
presents. The reader is assumed to be comfortable with linear algebra,
multivariable calculus, and the basics of machine learning; the geometric and
group-theoretic prerequisites are developed as needed.
